\lecture{44}{Сходимость FedAvg}

\sourcevideo
    {https://www.youtube.com/playlist?list=PLOzRYVm0a65fhKDS8cYIC7YCuVGJ4FOJx}
    {Week 11: Lecture 44: Convergence Analysis of FL}

Рассматривается задача
\[
    \min_{x\in\mathbb R^d}F(x),
    \qquad
    F(x)=\sum_{i=1}^{K}p_iF_i(x),
\]
где
\[
    p_i\ge0,
    \qquad
    \sum_{i=1}^{K}p_i=1.
\]
В FedAvg на каждом коммуникационном раунде выбирается часть клиентов,
каждый из них выполняет несколько локальных стохастических шагов, а
сервер усредняет полученные модели. Сходимость определяется совместным
действием стохастического шума, неполного участия клиентов и дрейфа
локальных траекторий.

\section{Предпосылки и динамика FedAvg}

Предполагается, что каждая локальная функция имеет
\(L\)-липшицев градиент:
\[
    \lVert\nabla F_i(x)-\nabla F_i(y)\rVert
    \le
    L\lVert x-y\rVert.
\]
Эквивалентная верхняя оценка имеет вид
\[
    F_i(y)
    \le
    F_i(x)
    +\langle\nabla F_i(x),y-x\rangle
    +\frac L2\lVert y-x\rVert^2.
\]

Клиент использует стохастический градиент \(g_i(x;\xi)\), для которого
\[
    \mathbb E_\xi[g_i(x;\xi)]
    =
    \nabla F_i(x)
\]
и
\[
    \mathbb E_\xi
    \left[
        \lVert
            g_i(x;\xi)-\nabla F_i(x)
        \rVert^2
    \right]
    \le
    \sigma^2.
\]
Следовательно,
\[
    \mathbb E_\xi\lVert g_i(x;\xi)\rVert^2
    \le
    \lVert\nabla F_i(x)\rVert^2+\sigma^2.
\]

Неоднородность локальных функций описывается условием
\[
    \sum_{i=1}^{K}
    p_i\lVert\nabla F_i(x)\rVert^2
    \le
    \beta^2\lVert\nabla F(x)\rVert^2
    +\kappa^2,
\]
где \(\beta^2\ge1\) и \(\kappa^2\ge0\). Параметр \(\beta\)
характеризует относительное различие направлений, а \(\kappa\) ---
остаточное различие локальных градиентов около стационарной точки
глобальной функции. В однородном идеализированном случае можно взять
\[
    \beta=1,
    \qquad
    \kappa=0.
\]

На раунде \(t\) выбранный клиент начинает из глобальной модели
\[
    x_i^{t,0}=x^t
\]
и выполняет \(\tau\) шагов
\[
    x_i^{t,s+1}
    =
    x_i^{t,s}
    -
    \eta g_i(x_i^{t,s};\xi_i^{t,s}),
    \qquad
    s=0,\ldots,\tau-1.
\]
Затем сервер агрегирует конечные локальные модели:
\[
    x^{t+1}
    =
    \sum_{i\in S_t}q_i^t x_i^{t,\tau},
    \qquad
    \sum_{i\in S_t}q_i^t=1.
\]

\section{Структура оценки сходимости}

При перечисленных предпосылках типичная невыпуклая оценка имеет
схематический вид
\[
\begin{aligned}
    \frac1T
    \sum_{t=0}^{T-1}
    \mathbb E
    \lVert\nabla F(x^t)\rVert^2
    \lesssim{}&
    \frac{F(x^0)-F_{\inf}}{\eta\tau T}
    \\
    &+
    \text{стохастическая ошибка}
    \\
    &+
    \text{ошибка локального дрейфа}.
\end{aligned}
\]
Первое слагаемое отвечает за начальную ошибку оптимизации.
Стохастическая часть зависит от дисперсии, размера мини-пакета и числа
участвующих клиентов. Дрейф возникает при \(\tau>1\), поскольку
градиенты вычисляются уже не в общей точке \(x^t\), а в различных
локальных точках \(x_i^{t,s}\).

Для анализа достаточно структуры оценки и её зависимости от параметров задачи; точные константы далее не используются.

Увеличение числа участников и размера мини-пакета обычно уменьшает
дисперсию агрегированного обновления. Увеличение числа локальных шагов
сокращает коммуникации, но усиливает накопленное отличие локальных
траекторий от глобального градиентного шага. Параметры \(\beta\) и
\(\kappa\) входят именно в оценку дрейфа и ухудшают её при росте
неоднородности.

\section{Участие клиентов и локальная работа}

Пусть в одном раунде участвует доля \(C\in(0,1]\) клиентов, так что
\[
    m\approx CK.
\]
При большем \(m\) агрегированное направление обычно имеет меньшую
дисперсию, поэтому может потребоваться меньше раундов. Одновременно
увеличиваются объём передачи, нагрузка на сервер и вероятность
ожидания медленного клиента.

Если клиент \(i\) имеет \(n_i\) объектов, использует мини-пакет
размера \(B\) и выполняет \(E\) локальных эпох, то число локальных
обновлений равно
\[
    \tau_i
    =
    E
    \left\lceil
        \frac{n_i}{B}
    \right\rceil.
\]
При фиксированном \(E\) увеличение \(B\) уменьшает число
последовательных шагов, но повышает объём обработки за один шаг и
требования к памяти. Поэтому реальное время зависит от оборудования и
параллелизма, а не только от числа итераций.

Увеличение \(E\) переносит больше работы внутрь коммуникационного
раунда. При умеренном значении это уменьшает число обменов с сервером.
При слишком большом \(E\) локальная модель начинает двигаться главным
образом к минимуму \(F_i\), а не к минимуму глобальной функции.

Время одного раунда складывается из рассылки модели, локального
обучения, обратной передачи и серверной агрегации. Параметры
\(\beta\) и \(\kappa\) непосредственно не задают длительность
локального шага, но влияют на число раундов, необходимое для достижения
заданной точности.

\section{Происхождение локального дрейфа}

Разворачивая локальную рекурсию, получаем
\[
    x_i^{t,\tau}
    =
    x^t
    -
    \eta
    \sum_{s=0}^{\tau-1}
    g_i(x_i^{t,s};\xi_i^{t,s}).
\]
После серверного усреднения
\[
    x^{t+1}
    =
    x^t
    -
    \eta
    \sum_{i\in S_t}q_i^t
    \sum_{s=0}^{\tau-1}
    g_i(x_i^{t,s};\xi_i^{t,s}).
\]
При \(\tau=1\) все локальные градиенты вычисляются в общей точке
\(x^t\). При \(\tau>1\) точки \(x_i^{t,s}\) начинают различаться,
поэтому усредняется сумма градиентов, вычисленных вдоль разных
траекторий.

Если локальные функции близки, то
\[
    \nabla F_i(x)
    \approx
    \nabla F(x),
\]
и несколько локальных шагов остаются хорошим приближением глобального
обновления. При сильной неоднородности направления
\(\nabla F_i(x)\) могут заметно различаться, а накопленный локальный
шаг становится смещённым относительно \(\nabla F(x)\). Поэтому
ошибка обычно растёт с \(\tau\), шагом \(\eta\) и параметрами
неоднородности.

\section{Точный пример с двумя клиентами}

Рассмотрим
\[
    F_1(x)=(x-1)^2,
    \qquad
    F_2(x)=2(x-5)^2
\]
и веса
\[
    p_1=p_2=\frac12.
\]
Тогда
\[
    F(x)
    =
    \frac12(x-1)^2+(x-5)^2
\]
и
\[
    F'(x)
    =
    (x-1)+2(x-5)
    =
    3x-11.
\]
Следовательно,
\[
    x^\star=\frac{11}{3}.
\]

Пусть оба клиента начинают из \(x^t\) и выполняют по одному точному
градиентному шагу:
\[
    x_1^{t,1}
    =
    x^t-2\eta(x^t-1),
\]
\[
    x_2^{t,1}
    =
    x^t-4\eta(x^t-5).
\]
После равномерного усреднения
\[
\begin{aligned}
    x^{t+1}
    &=
    \frac12x_1^{t,1}
    +
    \frac12x_2^{t,1}
    \\
    &=
    x^t
    -
    \eta\bigl[(x^t-1)+2(x^t-5)\bigr]
    \\
    &=
    x^t-\eta F'(x^t).
\end{aligned}
\]
При одном локальном шаге FedAvg здесь точно совпадает с градиентным
спуском по глобальной функции.

Если же каждый клиент полностью минимизирует собственную функцию, то
\[
    x_1^\star=1,
    \qquad
    x_2^\star=5,
\]
и сервер получает
\[
    \frac{x_1^\star+x_2^\star}{2}=3.
\]
Однако
\[
    3\ne\frac{11}{3}.
\]
Таким образом, если \(x_i^\star\in\operatorname*{arg\,min}F_i\),
то в общем случае
\[
    \sum_i p_i x_i^\star
    \ne
    x^\star,
    \qquad
    x^\star
    \in
    \operatorname*{arg\,min}_x
    \sum_i p_iF_i(x).
\]
Пример точно показывает источник смещения при чрезмерной локальной
оптимизации.

\section{Компромисс и границы анализа}

Увеличение \(m\) улучшает оценку глобального направления, но делает
каждый раунд дороже. Увеличение \(E\) и \(\tau\) уменьшает частоту
коммуникаций, но усиливает дрейф. Уменьшение \(B\) увеличивает число
локальных шагов и стохастический шум. Поэтому параметры должны
выбираться совместно с учётом стоимости связи, мощности клиентов,
неоднородности данных и требуемой точности.

Рассмотренная схема предполагает гладкость локальных функций,
несмещённые стохастические градиенты, ограниченную дисперсию и условие
ограниченной несхожести. Она не описывает вычислительную
неоднородность устройств, выпадение клиентов, асинхронность, задержки,
защищённую агрегацию и ограничения конфиденциальности. Эти эффекты
изменяют как время одного раунда, так и саму динамику алгоритма.
